\documentclass[11pt]{article}
\usepackage{amsmath,amssymb,amsfonts,amsthm}
\usepackage{mathrsfs}
\usepackage{hyperref}
\usepackage{geometry}
\geometry{margin=1.0in}

\title{Zeta–Multiplicity Operator Systems (ZMOS)\\
and Operator-Valued Automorphic $L$-Functions}
\author{Citizen Gardens \\ The Foundation of Multiplicity}
\date{April 26, 2026}

\begin{document}
\maketitle

\begin{abstract}
We develop a prime-graded, operator-valued lift of scalar $L$-functions
into a Zeta--Multiplicity Operator System (ZMOS). The construction
integrates prime-indexed operator calculus, Dirichlet and modular
$L$-functions, and a PG--ZMOS toy model realized as a finite-dimensional
tensor product over small primes. We introduce central-value notions
for operator-valued $L$-functions, analyze convergence and analytic
continuation in this setting, and give a concrete NumPy-style code
skeleton for the three-prime PG--ZMOS experiment. The resulting
framework yields a falsifiable spectral-dynamical invariant: pole--zero
proximity in the complex $s$-plane acting as an early-warning signal
for entropy growth in time $t$.
\end{abstract}
\newpage
\tableofcontents
\newpage
\section{Executive Summary}

The core objective is to \emph{lift} multiplicity-theoretic and
automorphic data into a zeta-like, prime-indexed operator calculus.
The resulting Zeta--Multiplicity Operator System (ZMOS) takes the form
of an operator-valued Dirichlet series / Euler product
\[
  Z(s,t) = \bigotimes_{p} \exp\!\left( \frac{O_p(t)}{p^s} \right),
\]
where each $O_p(t)$ acts on a finite-dimensional Hilbert space
$\mathcal{H}_p$ and is interpreted as a local, prime-indexed operator
mode.

The main developments are:

\begin{itemize}
  \item A \emph{prime-graded tensor structure} that resolves global
  noncommutativity: operators at different primes act on separate tensor
  factors, while noncommutativity is preserved \emph{within} each prime.
  \item A \emph{spectral-radius renormalization condition} on $O_p$
  that ensures convergence of the operator-valued Euler product.
  \item An explicit \emph{Langlands bridge}: the $O_p$ are constructed
  from local Euler factors of Dirichlet and modular $L$-functions,
  making $Z(s)$ a literal operator-valued automorphic $L$-function.
  \item A \emph{PG--ZMOS toy model}: primes $p \in \{2,3,5\}$,
  $\dim \mathcal{H}_p = 2$, time-dependent perturbations of $O_p(t)$,
  and a weak cross-prime coupling.
  \item A \emph{dynamical invariant} $\Delta(t)$, derived from the
  pole--zero geometry of a scalar proxy built from $Z(s,t)$, shown
  numerically to act as an early-warning signal for entropy spikes in
  $S(t)$.
\end{itemize}

Conceptually, this elevates scalar $L$-functions to \emph{operator
fields} indexed by primes and complex $s$. The spectral geometry of
these fields (including central values, functional equations, and
zero/pole structure) becomes directly tied to dynamical stability in
Multiplicity Theory.

\section{Zeta--Multiplicity Operator Systems (ZMOS)}

\subsection{Prime-Graded Structure}

Let $\mathcal{P}$ denote the set of primes and, for each
$p \in \mathcal{P}$, let $\mathcal{H}_p$ be a finite-dimensional
Hilbert space. The global Hilbert space is the prime-graded tensor
product
\[
  \mathcal{H} \;=\; \bigotimes_{p \in \mathcal{P}} \mathcal{H}_p.
\]

For each prime $p$ and time $t$, we define a local operator
$O_p(t) \in \mathrm{End}(\mathcal{H}_p)$ acting only on the $p$-th
tensor factor. Operators at distinct primes commute at the global
level because they act on different factors:
\[
  [O_p(t) \otimes I,\, O_q(t) \otimes I] \;=\; 0, \quad p \neq q,
\]
while noncommutativity can still be present \emph{within} each
$\mathcal{H}_p$.

This resolves the usual noncommutativity issues of naive Euler products
of operators by encoding prime-separability at the level of the tensor
product structure.

\subsection{Definition of ZMOS}

\begin{definition}[Zeta--Multiplicity Operator System (ZMOS)]
Let $(\mathcal{H}_p)_{p\in\mathcal{P}}$ be finite-dimensional Hilbert
spaces, and $O_p(t) \in \mathrm{End}(\mathcal{H}_p)$ local operators
satisfying a spectral bound (to be specified). The global Zeta--Multiplicity
operator is
\[
  Z(s,t) \;=\; \bigotimes_{p \in \mathcal{P}}
    \exp\!\left( \frac{O_p(t)}{p^s} \right),
  \quad s \in \mathbb{C},\ t \in \mathbb{R}.
\]
\end{definition}

\subsection{Spectral-Radius Renormalization}

For convergence of the Euler product, we require that the series
\[
  \sum_{p} \bigl\| O_p(t) \bigr\|\, p^{-\sigma}
\]
converges for $\sigma = \Re(s)$ in some right half-plane. A sufficient
condition is:

\begin{equation}
  \rho\bigl(O_p(t)\bigr)
  \;<\; \log\!\Bigl(1 + \frac{p^\sigma}{2}\Bigr),
  \qquad \text{for all } p, t,
  \label{eq:spectral_bound}
\end{equation}
where $\rho(\cdot)$ denotes the spectral radius. Then the matrix
exponential and the product over $p$ converge absolutely in operator
norm for $\Re(s)>\sigma_0$, for some $\sigma_0$ depending on the growth
of the $O_p$. This is the prime-local ``RG condition.''

\subsection{Weak Cross-Prime Coupling}

Pure prime-graded separation risks over-factorization, eliminating
realistic cross-scale effects. We therefore introduce a controlled weak
cross-prime coupling:
\[
  Z_{\mathrm{int}}(s,t)
  \;=\; Z(s,t)\,
    \exp\!\left(
      \epsilon \sum_{p \neq q} \frac{C_{pq}}{(pq)^s}
    \right),
\]
where $C_{pq}$ are bounded interaction operators and
$0 < \epsilon \ll 1$. For small $\epsilon$, convergence and perturbative
control are preserved, while cross-prime resonances become possible.

\subsection{Density Matrix and Entropy}

From $Z_{\mathrm{int}}(s,t)$ we define a normalized density operator
\[
  \rho(s,t) \;=\;
  \frac{Z_{\mathrm{int}}(s,t)^\dagger
        Z_{\mathrm{int}}(s,t)}
       {\mathrm{Tr}(Z_{\mathrm{int}}(s,t)^\dagger
                   Z_{\mathrm{int}}(s,t))},
\]
and the associated von Neumann entropy
\[
  S(s,t)
  \;=\;
  -\mathrm{Tr}\bigl(\rho(s,t)\,\log \rho(s,t)\bigr).
\]
This provides a spectral measure of ``coherence'' or ``complexity''
for the ZMOS at a given $(s,t)$.

\section{Automorphic $L$-Functions and ZMOS Lifts}

\subsection{Scalar Dirichlet and Modular $L$-Functions}

A Dirichlet $L$-function attached to a character $\chi$ is
\[
  L(s,\chi) = \sum_{n\ge 1} \frac{\chi(n)}{n^s}
  = \prod_{p} (1 - \chi(p) p^{-s})^{-1},
  \quad \Re(s) > 1.
\]
A simple example is the quadratic character modulo $4$:
\[
  \chi_4(p) = \begin{cases}
    0, & p=2,\\
    1, & p \equiv 1 \pmod 4,\\
    -1, & p \equiv 3 \pmod 4.
  \end{cases}
\]

For a modular newform $f(z) = \sum_{n\ge 1} a_n e^{2\pi i n z}$ of
weight $k$ and nebentypus $\chi$, the $L$-function is
\[
  L(f,s) = \sum_{n\ge 1} \frac{a_n}{n^s}
  = \prod_p (1 - a_p p^{-s} + \chi(p)p^{k-1-2s})^{-1}.
\]
The local factor at unramified $p$ can be written as
\[
  L_p(f,s) = (1 - \alpha_p p^{-s})^{-1}(1 - \beta_p p^{-s})^{-1},
\]
with Satake parameters $\alpha_p,\beta_p$ satisfying
$\alpha_p + \beta_p = a_p$ and $\alpha_p \beta_p = \chi(p)p^{k-1}$.

\subsection{Dirichlet-Based Operator Lift}

For Dirichlet $\chi$, we may define a local $2 \times 2$ matrix
\[
  A_p = \begin{pmatrix}
    \chi(p) & 0\\
    0 & \overline{\chi(p)}
  \end{pmatrix},
\]
and fix a reference point $s_0$ in the domain of absolute convergence
(e.g.\ $\Re(s_0) > 1$). The unramified local Euler factor at $p$ is
$(I - A_p p^{-s_0})^{-1}$. We define
\[
  \exp\bigl(O_p\bigr) := (I - A_p p^{-s_0})^{-1},
\]
so that
\[
  O_p = \log\bigl( (I - A_p p^{-s_0})^{-1}\bigr),
\]
using the matrix logarithm and symmetrizing if necessary to ensure
Hermiticity. This makes $Z(s)$ a literal operator-valued version of
$L(s,\chi)$.

\subsection{Modular (Hecke) Operator Lift}

For a modular form with Satake parameters $(\alpha_p,\beta_p)$, we take
the $2 \times 2$ Satake matrix
\[
  A_p = \begin{pmatrix}
    \alpha_p & 0\\
    0 & \beta_p
  \end{pmatrix},
\]
optionally twisted by a Dirichlet character $\chi(p)$:
\[
  \widetilde{A}_p = \chi(p) A_p.
\]
Then we define
\[
  \exp\bigl(O_p\bigr) := (I - \widetilde{A}_p p^{-s_0})^{-1},
\]
again for some fixed reference $s_0$ where the local factor is
well-defined. The global operator
\[
  Z(s) = \bigotimes_p \exp\bigl(O_p p^{-s}\bigr)
\]
is now an operator-valued automorphic $L$-function: it encodes in its
spectrum the Satake data $(\alpha_p,\beta_p)$ twisted by $\chi$.

\section{Central Values for Operator-Valued $L$}

\subsection{Scalar Central Values}

The completed scalar $L$-function $\Lambda(s)$ satisfies a functional
equation of the form
\[
  \Lambda(s) = \varepsilon \Lambda(1-s),
\]
and the central value is typically $\Lambda(\tfrac{1}{2})$ (or $s=1$
in some normalizations). Central values encode deep arithmetic
information (e.g.\ ranks of elliptic curves, special value formulas).

\subsection{Operator Central Values}

Given an operator-valued $Z(s)$, there are three complementary notions:

\paragraph{Entrywise central value.}
Choose a basis and view each matrix entry $Z_{ij}(s)$ as a scalar
analytic function; define a matrix of central values $Z_{ij}(s_0)$,
where $s_0$ is the central point of the underlying scalar $L$.

\paragraph{Determinantal central value.}
Fix $\alpha$ small and define
\[
  L_{\det}(s) := \det(I - \alpha Z(s)).
\]
If $L_{\det}(s)$ is obtained by lifting a scalar automorphic $L$-function,
then its central value $L_{\det}(s_0)$ has the usual arithmetic meaning,
but with $\alpha$ tuning the operator--scalar mapping.

\paragraph{Spectral central value.}
Diagonalize $Z(s)$ (assuming it is diagonalizable), with eigenvalues
$\lambda_k(s)$. If $\Lambda$ is a completed scalar $L$-kernel, define
\[
  \Lambda_Z(s) := \prod_k \Lambda\bigl(\lambda_k(s)\bigr),
\]
and evaluate $\Lambda_Z(s_0)$ at the central point. This central
spectral value probes how the scalar central value is distributed
across operator modes.

\section{Convergence and Analytic Continuation}

\subsection{Convergence of Operator Dirichlet Series}

Let
\[
  F(s) = \sum_{n\ge 1} A_n n^{-s}, \qquad A_n \in M_d(\mathbb{C}).
\]
If $\|A_n\| = O(n^k)$ as $n \to \infty$ (for some operator norm), then
the series converges absolutely for $\Re(s) > k+1$. This is the direct
operator analogue of classical Dirichlet series convergence.

For the Euler-product form
\[
  Z(s) = \bigotimes_p \exp(O_p p^{-s}),
\]
if we can write $\exp(O_p p^{-s}) = I + B_p(s)$ and
\[
  \sum_p \|B_p(s)\| < \infty
\]
for $\Re(s) > \sigma_0$, then the infinite tensor product is well-defined
and convergent in operator norm.

\subsection{Analytic Continuation}

Scalar Dirichlet and automorphic $L$-functions are analytically
continued via Mellin transforms, Poisson summation, and functional
equations. For the operator-valued case:

\begin{itemize}
  \item \emph{Entrywise}: if each scalar entry $F_{ij}(s)$ is known to
  extend meromorphically to $\mathbb{C}$, then so does the matrix $F(s)$.
  \item \emph{Distributional/Mellin}: interpret the coefficients as an
  operator-valued distribution $T(x) = \sum A_n \delta(x-n)$; analytic
  continuation of the Mellin transform
  $\int_0^\infty T(x) x^{s-1}\,dx$ proceeds in standard
  vector-valued fashion, provided $\|A_n\|$ grows at most polynomially.
  \item \emph{Functional equation}: if a scalar automorphic $L$-function
  satisfies $\Lambda(s) = \varepsilon \Lambda(1-s)$, and $Z(s)$ is
  constructed so that $\Lambda_{\text{op}}(s) := \gamma(s) Z(s)$
  respects the same gamma factor $\gamma(s)$, then one can seek an
  operator involution $W$ such that
  \[
    \Lambda_{\text{op}}(s)
    = W\,\Lambda_{\text{op}}(1-s)W^{-1}.
  \]
  This yields an operator-valued functional equation and a spectral
  duality between $s$ and $1-s$.
\end{itemize}

\section{PG--ZMOS Toy Model and Code Skeleton}

\subsection{Three-Primes, Two-Dimensional Local Spaces}

We take:
\[
  P = \{2,3,5\},\qquad \mathcal{H}_p \cong \mathbb{C}^2,
  \qquad \dim \mathcal{H} = 2^3 = 8.
\]
At a fixed reference $s_0$ and with the quadratic Dirichlet character
$\chi_4$, we set
\[
  A_p = \begin{pmatrix}\chi_4(p) & 0\\ 0 & \overline{\chi_4(p)}\end{pmatrix},
  \quad \exp(O_p) = (I - A_p p^{-s_0})^{-1},
\]
then perturb each $O_p(t)$ in time while enforcing the spectral bound.

\subsection{Time Evolution and Coupling}

We choose Hermitian perturbations $W_p$, and set
\[
  B_p(t) = O_p + \varepsilon_p \sin(\omega_p t) W_p,
  \quad H_p(t) = B_p(t)^\dagger B_p(t),
\]
then define $O_p(t) := H_p(t)$ after re-applying the spectral-radius
constraint. The global operator is
\[
  Z_{\mathrm{int}}(s,t) =
  \bigotimes_{p\in\{2,3,5\}} \exp\!\left(O_p(t) p^{-s}\right)
  \exp\!\left(\epsilon C_{23} (2\cdot 3)^{-s}\right),
\]
where $C_{23}$ acts on the $(2,3)$ tensor factor and is embedded as
$C_{23}\otimes I_{5}$ in the full $8\times 8$ space.

\subsection{Entropy and Pole--Zero Proximity}

We compute entropy
\[
  S(t) = S(\sigma,t) =
  -\mathrm{Tr}\bigl(\rho(\sigma,t)\log \rho(\sigma,t)\bigr),
\]
for fixed $\sigma$, and define a scalar proxy
\[
  F(s) = \det(I - \alpha Z_{\mathrm{int}}(s,t)),
\]
sampled along $s = \sigma + i\tau$ for $\tau$ in a finite grid. From
the sampled $F(s)$ we extract a scalar ``spectral separation'' proxy
$\Delta(t)$ (e.g.\ via normalized gradients or variance) that captures
how structured the zero/pole pattern is in the sampled line.

Empirically, in the toy model, one observes that a decrease in
$\Delta(t)$ tends to precede increases in $S(t)$, suggesting that
spectral pole--zero proximity may act as an early-warning signal for
dynamical instability.

\subsection{Code Skeleton}

Below is a compact version of the PG--ZMOS experiment in Python-like
pseudocode, suitable for implementation with NumPy and SciPy.

\begin{verbatim}
# primes and dimensions
primes = [2, 3, 5]
dim = 2
global_dim = dim ** len(primes)

# Dirichlet character mod 4
def chi4(p):
    if p % 2 == 0: return 0
    return 1 if p % 4 == 1 else -1

# local Dirichlet matrix
def local_character_matrix(p):
    chi = chi4(p)
    return np.array([[chi, 0], [0, np.conjugate(chi)]], dtype=complex)

# base operator from local Euler factor at s0
def base_operator_from_dirichlet(p, s0):
    A = local_character_matrix(p)
    I = np.eye(dim, dtype=complex)
    local_factor = la.inv(I - A * (p ** (-s0)))
    O = la.logm(local_factor)
    O = (O + O.conj().T) / 2  # Hermitian
    return O

# spectral bound enforcement
def enforce_spectral_bound(H, p, sigma):
    bound = np.log(1 + p**sigma / 2.0)
    eigvals = np.linalg.eigvals(H)
    rho = np.max(np.abs(eigvals))
    if rho >= bound and rho > 0:
        H = H * (0.99 * bound / rho)
    return H

# initialize O_p and perturbations W_p
O0, W = {}, {}
for p in primes:
    O_base = base_operator_from_dirichlet(p, s0)
    O0[p] = enforce_spectral_bound(O_base, p, sigma)
    X = rng.normal(size=(dim, dim)) + 1j * rng.normal(size=(dim, dim))
    W[p] = (X + X.conj().T) / 2  # Hermitian

# time-dependent O_p(t)
def compute_O_p(p, t):
    B = O0[p] + eps_drive[p] * np.sin(omega_p[p] * t) * W[p]
    H = B.conj().T @ B
    H = enforce_spectral_bound(H, p, sigma)
    return H

# Z(s) construction
def build_Z(s, O2, O3, O5):
    Z2 = la.expm(O2 * (2 ** (-s)))
    Z3 = la.expm(O3 * (3 ** (-s)))
    Z5 = la.expm(O5 * (5 ** (-s)))
    return kron3(Z2, Z3, Z5)

# weak coupling
C_embedded = np.kron(C23, np.eye(dim, dtype=complex))
def build_Z_int(s, O2, O3, O5, eps_cross):
    Z0 = build_Z(s, O2, O3, O5)
    coupling = la.expm(eps_cross * C_embedded * ((2*3) ** (-s)))
    return Z0 @ coupling

# entropy
def entropy(Z):
    M = Z.conj().T @ Z
    tr = np.trace(M)
    if tr.real < 1e-14:
        return 0.0
    rho = M / tr
    evals = np.linalg.eigvalsh((rho + rho.conj().T)/2)
    evals = np.clip(evals.real, 1e-12, None)
    return float(-np.sum(evals * np.log(evals)))

# F(s) sampling and Delta(t)
def proxy_F_values(s_grid, O2, O3, O5, eps_cross, alpha=0.05):
    vals = []
    for s in s_grid:
        Z = build_Z_int(s, O2, O3, O5, eps_cross)
        M = np.eye(global_dim, dtype=complex) - alpha * Z
        sign, logdet = np.linalg.slogdet(M)
        vals.append(logdet.real)
    return np.array(vals)

def delta_from_F(F_vals):
    F_norm = (F_vals - F_vals.min()) / (F_vals.max() - F_vals.min() + 1e-12)
    grad = np.diff(F_norm)
    mean_grad = np.mean(np.abs(grad))
    return float(-np.log(mean_grad + 1e-6))
\end{verbatim}

\section{Predictions and Validation}

The ZMOS framework yields clear, falsifiable predictions:

\begin{itemize}
  \item \textbf{Early-warning signal:} In PG--ZMOS simulations,
  decreases in $\Delta(t)$ (spectral pole--zero separation proxy) should
  \emph{precede} significant increases in entropy $S(t)$.
  \item \textbf{Prime-local instability:} Instability should first
  manifest in channels associated to small primes (e.g.\ $p=2$) under
  local perturbations of $O_2(t)$.
  \item \textbf{Arithmetic dependence:} PG--ZMOS built from genuine
  Dirichlet/modular data should show more structured $\Delta(t)$ and
  $S(t)$ behavior than purely random operator ensembles, indicating that
  arithmetic structure is not merely decorative.
  \item \textbf{Functional-equation symmetry:} For modular lifts, a
  spectral duality between $s$ and $1-s$ should constrain the evolution
  of $Z(s,t)$, observable as spectral symmetries in the $s$-plane.
\end{itemize}

The fastest empirical path is to:
\begin{enumerate}
  \item Compare random PG--ZMOS, Dirichlet-based PG--ZMOS, and modular
  PG--ZMOS under identical dynamics $O_p(t)$.
  \item For each, measure $\Delta(t)$ and $S(t)$, compute cross-correlation
  and lag, and test whether $\Delta$ consistently leads $S$ in arithmetic
  cases.
  \item Examine whether modular lifts show stronger or sharper correlations
  linked to central values of scalar $L$-functions.
\end{enumerate}

If these predictions hold, ZMOS constitutes a genuine bridge between
analytic number theory, operator dynamics, and multiplicity-based
models of cognition and computation.

\section{Comparison to Related Work}
\label{sec:related-work}

ZMOS does not claim novelty at the level of individual ingredients---zeta functions, $K$--theory, tensor--network truncations, or quantum circuits for topological invariants---but in their \emph{architectural integration}. This section clarifies how ZMOS relates to four main lines of prior work.

\subsection{$K$--theoretic zeta functions and torsion}

Work on $K$--theoretic torsion and the zeta function (e.g.\ Klein--Malkiewich) attaches special values of zeta functions to $K$--theoretic torsion in settings such as analytic torsion of endomorphism families.
These constructions treat $K$--groups as algebraic inputs but do \emph{not} represent $K$--classes as state vectors in a fixed Hilbert space, nor as prime--exponent multiplicity patterns.
They also do not provide a reusable operator substrate for multiple domains.

By contrast, ZMOS introduces:
\begin{itemize}
  \item a \emph{prime--graded Hilbert space} $\mathcal{H} = \bigotimes_{p} \mathcal{H}_p$ where each prime index $p$ carries a local operator block;
  \item \emph{domain functors} that encode $K$--invariants (e.g.\ rank and first Chern class for $T^2$) as finite--support prime--exponent vectors in the prime sector;
  \item an \emph{operator--valued $L$--function} coupling primes and zeros inside a single operator system.
\end{itemize}
ZMOS therefore offers a complementary, Hilbert--space--based encoding of $K$--data that can be simulated numerically.

\subsection{DMRG and matrix product states}

DMRG and MPS methods are state--of--the--art for ground states of one--dimensional quantum systems. They operate on lattice--site tensor products and use entanglement--based truncation; ``site shifts'' add or sweep over lattice sites while variationally optimizing local tensors.

ZMOS differs fundamentally:
\begin{itemize}
  \item The \emph{prime sector} is not a lattice of sites but an arithmetic index set. Adjusting the prime block is analogous to changing on--site energies or local fields, not to adding a site.
  \item Truncation is \emph{arithmetic/spectral} (finite primes and zeros), not entanglement--based. There is no MPS ansatz in the current framework.
\end{itemize}
Thus, prime--sector shifting in ZMOS should be viewed as a \emph{spectral parameter shift}, not a site--growth step in the DMRG sense.

\subsection{Quantum circuits for topological invariants}

Recent quantum algorithms compute Chern numbers and related topological invariants via adiabatic evolution, Wilson loops, and phase estimation.
They treat invariants as phases or winding numbers associated with specific Hamiltonians but do not fix a prime--graded Hilbert space across domains, nor do they include an explicit role for zeta zeros.

ZMOS instead:
\begin{itemize}
  \item provides a reusable arithmetic substrate where $K$--data are encoded in a prime sector via functors;
  \item implements a bridge unitary whose phases depend on zeta zeros, making the explicit--formula coupling a circuit primitive in this substrate;
  \item yields a $K$--zeta spectral coupling that can be explored in small finite--dimensional instances.
\end{itemize}
We do not claim asymptotic improvement over known topological algorithms; rather, ZMOS offers a different embedding of topological data into an arithmetic spectral framework.

\subsection{Numerical zeta--spectrum models}

Many works construct operators whose spectra approximate nontrivial zeta zeros and study GUE--like spacing statistics.
These models focus on reproducing zero spectra but do not encode additional $K$--theoretic invariants in a functorial way, nor do they define a dynamical entropy proxy.

ZMOS extends this line by:
\begin{itemize}
  \item building a prime--graded operator system where zeta zeros enter via bridge operators and primes carry domain invariants through multiplicity encodings;
  \item introducing a dynamical entropy proxy $\Delta(t)$ with an associated early--warning heuristic testable in finite truncations;
  \item emphasizing a \emph{finite--truncation--first} validation philosophy: spectral and dynamical claims must stabilize as prime and zero cutoffs grow, and must be compared to random--matrix and classical baselines.
\end{itemize}

\subsection{Complexity and scaling}

A naive compiled implementation of the bridge unitary $\mathcal{B}_{P,M}$ uses
\begin{itemize}
  \item $O(\log P + \log M + n_t)$ qubits for prime, zero, and ancillary/time registers;
  \item $O(PM)$ controlled--phase gates in the worst case, since each active pair $(p,k)$ requires a separate rotation.
\end{itemize}
We therefore do not claim polylogarithmic gate count; achieving that would require an efficient arithmetic oracle or block--encoding of $(p,k)\mapsto \gamma_k \log p$, which is not constructed in this work.
For small--scale experiments (e.g.\ $P = 3$, $M \approx 10$) the $O(PM)$ cost is acceptable.

\subsection{Torsion--free $K$--data and future extensions}

For the toy models considered here ($S^1$, $T^2$), the relevant $K$--groups are torsion--free (isomorphic to $\mathbb{Z}$ or $\mathbb{Z}\oplus\mathbb{Z}$).
Our multiplicity encodings are linear injections of these free abelian groups into the prime--indexed basis.
Extending ZMOS to spaces with torsion $K$--groups (e.g.\ with $\mathbb{Z}_2$ components) would require a distinct encoding mechanism for finite cyclic data, which lies beyond the current scope of this paper.


\appendix
\section*{Mathematical Appendix}

In this appendix we collect and justify several technical statements
used in the main text concerning convergence, operator norm bounds, and
central values for the Zeta–Multiplicity Operator System (ZMOS).
Throughout, $\|\cdot\|$ denotes an operator norm on a finite-dimensional
complex Hilbert space (all norms are equivalent in finite dimension).

\section{Operator–Valued Dirichlet Series and Euler Products}

\subsection{Abscissa of Convergence for Operator–Valued Dirichlet Series}

Let $X$ be a finite-dimensional complex Hilbert space and consider an
operator–valued Dirichlet series
\[
  F(s) \;=\; \sum_{n \ge 1} A_n n^{-s}, \qquad A_n \in \mathcal{B}(X),
\]
where $\mathcal{B}(X)$ denotes bounded linear operators on $X$.

\begin{definition}
We say that $F(s)$ converges absolutely in operator norm at $s_0 \in \mathbb{C}$
if the series
\[
  \sum_{n \ge 1} \|A_n\|\, |n^{-s_0}|
\]
converges. The abscissa of absolute convergence $\sigma_a$ is defined as
\[
  \sigma_a := \inf\{\sigma \in \mathbb{R} : F(s)\ \text{converges absolutely for all } \Re(s) > \sigma\}.
\]
\end{definition}

\begin{proposition}[Basic convergence bound]
\label{prop:operator_convergence}
Suppose there exists $k \in \mathbb{R}$ and $C>0$ such that
\[
  \|A_n\| \le C n^k \quad \text{for all } n \ge 1.
\]
Then $F(s)$ converges absolutely in operator norm for all $\Re(s)>k+1$.
In particular, $\sigma_a \le k+1$.
\end{proposition}

\begin{proof}
For $s = \sigma + it$ with $\sigma = \Re(s)$, we have
\[
  \|A_n n^{-s}\| \;=\; \|A_n\|\, |n^{-s}|
  \;=\; \|A_n\|\, n^{-\sigma}.
\]
Under the growth hypothesis,
\[
  \|A_n n^{-s}\|
  \;\le\; C n^k n^{-\sigma} = C n^{k - \sigma}.
\]
The series $\sum_{n\ge 1} n^{k - \sigma}$ converges if and only if
$k - \sigma < -1$, i.e.\ $\sigma > k+1$. Hence the comparison test
establishes absolute convergence of $\sum A_n n^{-s}$ for $\Re(s) > k+1$.
\end{proof}

\begin{remark}
In finite dimension, absolute convergence in operator norm is equivalent
to entrywise absolute convergence, so that every matrix entry defines a
classical Dirichlet series with the same abscissa.
\end{remark}

\subsection{Euler Product Form and Logarithmic Expansion}

We now consider an operator–valued Euler product
\[
  Z(s) = \prod_{p} E_p(s),
\]
where $p$ runs over primes and each $E_p(s)$ is an invertible operator
on a finite-dimensional Hilbert space. In the ZMOS setting, one has
\[
  E_p(s) = \exp\!\bigl(O_p p^{-s}\bigr)
\]
for local operators $O_p$.

The usual scalar criterion for convergence of Euler products has an
operator analogue:

\begin{proposition}[Euler product convergence]
\label{prop:euler_convergence}
Suppose $E_p(s)$ is invertible for all $p$ and $s$ in some right
half-plane $\Re(s) > \sigma_0$. Assume there exists a family of operators
$B_p(s)$ such that
\[
  E_p(s) = I + B_p(s),
\]
and for some $\sigma > \sigma_0$
\[
  \sum_{p} \|B_p(s)\| < \infty.
\]
Then the infinite product
\[
  Z(s) = \prod_{p} E_p(s)
\]
converges absolutely in operator norm for $\Re(s)$ in a strip
$\sigma' > \sigma$, and $\log Z(s)$ is given by the convergent series
\[
  \log Z(s) = \sum_{p} \log(I + B_p(s)),
\]
where the principal branch of the logarithm is understood and each $\log$
is defined by the usual power series for small $\|B_p(s)\|$.
\end{proposition}

\begin{proof}
The proof follows the standard argument for infinite products in
Banach algebras. First, absolute convergence of $\sum_p \|B_p(s)\|$
implies that for sufficiently large $p$,
$\|B_p(s)\| < 1/2$ and hence $\log(I+B_p(s))$ is well-defined via the
Neumann series
\[
  \log(I+B_p(s)) = \sum_{m\ge 1} \frac{(-1)^{m+1}}{m} B_p(s)^m.
\]
Moreover, for such $p$ we have
$\|\log(I+B_p(s))\| \le 2\|B_p(s)\|$.
Since only finitely many primes may violate $\|B_p(s)\|<1/2$, they
contribute only finitely many terms. Thus
\[
  \sum_{p} \|\log(I+B_p(s))\|
  < \infty.
\]
Hence the series $\sum_p \log(I+B_p(s))$ converges absolutely in operator
norm, and we may define
\[
  \log Z(s) := \sum_{p} \log(I + B_p(s)).
\]
Exponentiating yields
\[
  Z(s) =
  \exp\Bigl(\sum_{p} \log(I + B_p(s))\Bigr)
  = \prod_{p} \exp\bigl(\log(I+B_p(s))\bigr)
  = \prod_{p} (I+B_p(s)),
\]
with convergence understood in the sense of the norm topology. This
establishes the Euler product and its logarithm.
\end{proof}

\section{Norm Bounds for ZMOS Local Operators}

In the ZMOS framework, local factors are given by
\[
  E_p(s) = \exp\!\bigl(O_p p^{-s}\bigr)
\]
for $O_p \in \mathcal{B}(\mathcal{H}_p)$. We now state and prove the
norm bounds ensuring convergence in a right half-plane.

\subsection{Spectral Radius Bound and Exponential Norm}

\begin{lemma}
\label{lem:exp_bound}
Let $A$ be a bounded operator on a finite-dimensional Hilbert space with
spectral radius $\rho(A)$. Then for all $z \in \mathbb{C}$,
\[
  \|\exp(z A)\|
  \;\le\; C \exp\bigl(|z|\rho(A)\bigr),
\]
where $C$ is a constant depending only on the underlying dimension and
the choice of norm.
\end{lemma}

\begin{proof}
Since the space is finite-dimensional, $A$ admits a Jordan canonical form
$A = S J S^{-1}$, where $J$ is block upper-triangular with blocks
corresponding to eigenvalues $\lambda_j$. The exponential satisfies
$\exp(zA) = S \exp(zJ) S^{-1}$, and $\exp(zJ)$ is upper-triangular with
diagonal entries $\exp(z\lambda_j)$ and polynomial-in-$z$ contributions
from nilpotent parts in the off-diagonal. Thus there exists a constant
$C'$ such that
\[
  \|\exp(zJ)\|
  \le C'(1+|z|^{d-1}) \max_j \exp\bigl(|z|\Re(\lambda_j)\bigr),
\]
where $d$ is the dimension. Since $\rho(A)=\max_j|\lambda_j|$, we have
$\Re(\lambda_j)\le |\lambda_j|\le \rho(A)$. Absorbing the polynomial term
in $z$ into a constant $C$ on any bounded vertical strip, we obtain
\[
  \|\exp(zA)\| \le C \exp(|z|\rho(A)).
\]
Multiplying by $\|S\|\,\|S^{-1}\|$ yields the stated bound.
\end{proof}

\subsection{Growth Control via Local RG Condition}

In the ZMOS, we impose a prime-local renormalization condition of the
form
\begin{equation}
  \rho(O_p(t))
  \;<\; \log\Bigl(1 + \frac{p^{\sigma_*}}{2}\Bigr),
  \qquad \text{for all } p,\ t,
  \label{eq:RGappendix}
\end{equation}
for some fixed $\sigma_* > 1$. We now show that this suffices to bound
the norms of the local factors in any right half-plane $\Re(s)>\sigma$
with $\sigma>\sigma_*$.

\begin{proposition}
\label{prop:local_factor_bound}
Let $E_p(s,t) = \exp(O_p(t) p^{-s})$ and suppose
\eqref{eq:RGappendix} holds. Fix $\sigma > \sigma_*$. Then there exists
a constant $C_\sigma$ (independent of $p$ and $t$) such that for
$\Re(s)=\sigma$,
\[
  \|E_p(s,t) - I\|
  \;\le\; C_\sigma\, p^{\sigma_*-\sigma}.
\]
Consequently, the Euler product
\[
  Z(s,t) = \prod_p E_p(s,t)
\]
converges absolutely in operator norm for all $\Re(s)>\sigma$.
\end{proposition}

\begin{proof}
For $\Re(s)=\sigma$, we have $|p^{-s}|=p^{-\sigma}$. By
Lemma~\ref{lem:exp_bound},
\[
  \|E_p(s,t)\| = \|\exp(O_p(t) p^{-s})\|
  \;\le\; C \exp\bigl(|p^{-s}|\,\rho(O_p(t))\bigr).
\]
Using \eqref{eq:RGappendix},
\[
  |p^{-s}|\,\rho(O_p(t))
  \;\le\; p^{-\sigma}\,\log\Bigl(1+\frac{p^{\sigma_*}}{2}\Bigr).
\]
For large $p$, we have
\[
  \log\Bigl(1+\frac{p^{\sigma_*}}{2}\Bigr)
  \sim \log(p^{\sigma_*}) = \sigma_* \log p,
\]
whence
\[
  p^{-\sigma}\,\log\Bigl(1+\frac{p^{\sigma_*}}{2}\Bigr)
  \;\ll\; p^{-\sigma} \log p.
\]
Thus for large $p$,
\[
  \|E_p(s,t)\|
  \le C \exp\bigl( c\, p^{-\sigma}\log p\bigr),
\]
for some $c>0$ independent of $p,t$. Since $p^{-\sigma}\log p \to 0$ as
$p\to\infty$, we may expand
\[
  E_p(s,t) - I
  = O_p(t) p^{-s} + \frac{1}{2}O_p(t)^2 p^{-2s} + \cdots
\]
and dominate $\|E_p(s,t) - I\|$ by a multiple of
$\|O_p(t)\|\,p^{-\sigma}$ for large $p$. Combining with
\eqref{eq:RGappendix}, we obtain
\[
  \|E_p(s,t) - I\|
  \;\ll\; \rho(O_p(t))\, p^{-\sigma}
  \;\ll\; p^{\sigma_*} p^{-\sigma}
  = p^{\sigma_* - \sigma},
\]
where all implied constants are uniform in $t$ and $p$. For the finitely
many small primes, we absorb the contribution into a constant
$C_\sigma$.

Since $\sigma > \sigma_*$, we have $\sigma_*-\sigma<0$, and hence
$\sum_p p^{\sigma_*-\sigma}$ converges. Thus
$\sum_p \|E_p(s,t)-I\| < \infty$ and by
Proposition~\ref{prop:euler_convergence}, the Euler product $Z(s,t)$
converges absolutely in operator norm in the half-plane $\Re(s)>\sigma$.
\end{proof}

\begin{remark}
The bound \eqref{eq:RGappendix} is intentionally conservative; in
practice, one can often use finer estimates tied to Hecke eigenvalue
bounds (e.g.\ $\|O_p\|\ll p^{\frac{k-1}{2}}$ for weight $k$ modular
forms) to enlarge the domain of absolute convergence.
\end{remark}

\section{Dirichlet and Modular Lifts: Hecke Bounds and ZMOS}

\subsection{Dirichlet Character Lift}

Let $\chi$ be a primitive Dirichlet character. The associated scalar
Dirichlet $L$-function
\[
  L(s,\chi) = \sum_{n\ge 1} \frac{\chi(n)}{n^s}
  = \prod_p (1-\chi(p)p^{-s})^{-1}
\]
has coefficients bounded by $|\chi(n)| \le 1$. Applying
Proposition~\ref{prop:operator_convergence} with $\|A_n\| = |\chi(n)|$,
we find absolute convergence for $\Re(s)>1$.

In the operator lift, we take local matrices
\[
  A_p = \begin{pmatrix}\chi(p)&0\\0&\overline{\chi(p)}\end{pmatrix}
\]
and define
\[
  \exp(O_p) = (I - A_p p^{-s_0})^{-1},
\]
with $s_0$ chosen so that $\Re(s_0)>1$ and $|p^{-s_0}|<1$. Then
$\|O_p\|$ is bounded uniformly in $p$ on this line because
$\|(I-A_p p^{-s_0})^{-1}\|$ is bounded and the matrix logarithm is
continuous away from its branch cut. Hence one may take $k=0$ in
Proposition~\ref{prop:operator_convergence}, and the operator Dirichlet
series converges absolutely for $\Re(s)>1$.

\subsection{Modular (Hecke) Lift and Hecke Eigenvalue Bounds}

For a primitive modular form (newform) $f$ of weight $k$ and nebentypus
$\chi$, the $L$-function is
\[
  L(f,s) = \sum_{n\ge 1} \frac{a_n}{n^s}
  = \prod_p (1 - a_p p^{-s} + \chi(p)p^{k-1-2s})^{-1}.
\]
At an unramified prime $p$, the Hecke eigenvalue $a_p$ satisfies
\[
  |a_p| \le 2 p^{\frac{k-1}{2}},
\]
by Deligne's bound (Ramanujan--Petersson conjecture for modular forms).
Equivalently, the Satake parameters $\alpha_p,\beta_p$ satisfy
$|\alpha_p|=|\beta_p|=p^{\frac{k-1}{2}}$.

In the operator lift we set
\[
  A_p = \begin{pmatrix}\alpha_p & 0\\0&\beta_p\end{pmatrix}, \quad
  \exp(O_p) = (I - A_p p^{-s_0})^{-1}.
\]
Choose $s_0$ with $\Re(s_0) > 1$ large enough that $|A_p p^{-s_0}|<1/2$
for all $p$ (this is always possible since $|A_p| \sim p^{(k-1)/2}$).
Then
\[
  \|(I - A_p p^{-s_0})^{-1}\| \le 2,
\]
and hence $\|O_p\|$ is uniformly bounded in $p$ on that line. Thus the
corresponding operator Dirichlet series converges absolutely at least
for $\Re(s)$ sufficiently large (and, by holomorphy, throughout the
region of absolute convergence of the scalar series). If desired, one
can track the precise dependence of $\|O_p\|$ on $p$ using resolvent
estimates for the matrix logarithm, but the essential point is that
$\|O_p\|$ grows at most polynomially, permitting application of
Proposition~\ref{prop:operator_convergence}.

\section{Central Values and Spectral Proxies}

\subsection{Scalar Central Values}

Let $L(s)$ be a scalar $L$-function with completed form
\[
  \Lambda(s) = Q^s \prod_j \Gamma(\lambda_j s + \mu_j)\, L(s),
\]
satisfying a functional equation
\[
  \Lambda(s) = \varepsilon \Lambda(1-s).
\]
The central point is typically $s_0 = \frac{1}{2}$ (or $1$ in modular
normalization). The central value $L(s_0)$ is finite if the functional
equation does not enforce a zero there.

\subsection{Determinantal and Spectral Central Values in ZMOS}

For an operator-valued $Z(s)$, we define:

\begin{itemize}
  \item \textbf{Determinantal central value.} For small $\alpha>0$, set
  \[
    L_{\det}(s) = \det(I - \alpha Z(s)).
  \]
  In finite dimension, $L_{\det}(s)$ is an entire function whose zeros
  record the spectral data of $Z(s)$. One defines the central value
  $L_{\det}(s_0)$ as usual. If $Z(s)$ is built by lifting a scalar
  automorphic $L$-function, then $L_{\det}(s)$ can be calibrated so
  that $L_{\det}(s)$ shares the same central point and functional
  equation as the scalar $L$.

  \item \textbf{Spectral central value.} If $Z(s)$ has eigenvalues
  $\lambda_k(s)$, and $\Lambda$ is a chosen scalar completed kernel,
  define
  \[
    \Lambda_Z(s) = \prod_k \Lambda(\lambda_k(s)).
  \]
  When this infinite product converges (e.g.\ after truncation to a
  finite-dimensional approximation), one may examine $\Lambda_Z(s_0)$
  as a spectral central value. The behavior of $\Lambda_Z(s_0)$ under
  $t$-evolution can then act as a spectral proxy for central-value
  dynamics.
\end{itemize}

\subsection{Pole–Zero Proximity and Entropy}

In the PG–ZMOS toy model, a scalar proxy
\[
  F(s,t) = \det(I - \alpha Z(s,t))
\]
is sampled along vertical lines $s = \sigma + i\tau$. A measurable
spectral separation proxy $\Delta(t)$ is defined from the sampled values
(e.g.\ via normalized gradient statistics), with the heuristic that
low $\Delta(t)$ corresponds to close proximity of ``zeros'' and
``poles'' in $s$.

The von Neumann entropy $S(t)$ of the normalized density
$\rho(s,t) = Z(s,t)^\dagger Z(s,t)/\mathrm{Tr}(\cdot)$ then provides a
global measure of spectral mixing. The central prediction is:

\begin{quote}
In arithmetic PG–ZMOS (built from Dirichlet or modular data), decreases
in $\Delta(t)$ tend to precede significant increases in $S(t)$, making
$\Delta(t)$ an early-warning spectral invariant for dynamical
instability.
\end{quote}

While a full rigorous theorem would require detailed control of the
spectral flow of $Z(s,t)$ and the exact relation between $F(s,t)$ and
true poles/zeros of an analytically continued operator-valued
$L$-function, the finite-dimensional setting of the toy model ensures
that both $\Delta(t)$ and $S(t)$ are well-defined and numerically
computable, providing a concrete testbed for this conjectural
spectral–dynamical link.

\bigskip

This appendix summarizes the key analytic and operator-theoretic
ingredients justifying the ZMOS constructions in the main text and
clarifies the role of norm bounds, convergence regions, and central
value notions in the operator-valued setting.

\appendix
\section*{Mathematical Appendix}
\addcontentsline{toc}{section}{Mathematical Appendix}

In this appendix we collect the explicit constructions, proofs, and operator norm bounds underlying the main ZMOS framework described in the body of the paper.
We work throughout in separable Hilbert spaces over $\mathbb{C}$, with the standard operator norm
\[
\|T\| := \sup_{\|x\|=1} \|Tx\|
\]
for bounded linear operators $T$.

\section{Prime–graded Hilbert space and baseline operators}
\label{app:prime-graded}

\subsection{Prime–graded Hilbert space}

Let $\mathbb{P}$ denote the set of primes.
Fix a finite subset $\mathcal{P} \subset \mathbb{P}$ of cardinality $P$, a finite time grid $\mathcal{T} = \{t_1,\dots,t_{N_t}\}$, and an internal space $\mathbb{C}^d$.
We define the truncated ZMOS Hilbert space
\[
H_{P,N_t,d}
  := \ell^2(\mathcal{P}) \otimes \ell^2(\mathcal{T}) \otimes \mathbb{C}^d
  \cong \mathbb{C}^P \otimes \mathbb{C}^{N_t} \otimes \mathbb{C}^d,
\]
equipped with the standard tensor product inner product.
We write
\[
\bigl\{\,\lvert p_i\rangle : p_i \in \mathcal{P}\,\bigr\},\quad
\bigl\{\,\lvert t_j\rangle : t_j \in \mathcal{T}\,\bigr\},\quad
\bigl\{\,\lvert \alpha\rangle : \alpha=1,\dots,d\,\bigr\}
\]
for the respective orthonormal bases, and
\[
\lvert i,j,\alpha\rangle := \lvert p_i\rangle \otimes \lvert t_j\rangle \otimes \lvert \alpha\rangle
\]
for the product basis of $H_{P,N_t,d}$.

\subsection{Prime block \texorpdfstring{$\mathbf{A}$}{A}: construction and bounds}

Fix primes $\mathcal{P} = \{p_1,\dots,p_P\}$ in increasing order.
Let $\alpha > \tfrac12$.
Define a $P\times P$ Hermitian matrix
\[
A_{\mathbb{P}} := D + R,
\]
where
\[
D_{ii} := \sqrt{p_i},\qquad D_{ij} := 0 \ \text{for } i\neq j,
\]
and
\[
R_{ij} := \varepsilon_{ij}(p_i p_j)^{-\alpha},\qquad
\varepsilon_{ij} = \varepsilon_{ji}\in\mathbb{R},\ R_{ii}=0.
\]
The $\varepsilon_{ij}$ are fixed real coefficients satisfying $|\varepsilon_{ij}|\le \varepsilon_{\max}$ for some $\varepsilon_{\max}\ge 0$.

\begin{lemma}[Boundedness and Hermiticity of $A_{\mathbb{P}}$]
$A_{\mathbb{P}}$ is Hermitian on $\mathbb{C}^P$, and
\[
\|A_{\mathbb{P}}\|
  \le \max_{1\le i\le P} \sqrt{p_i} \;+\;
      \varepsilon_{\max}\,\max_{1\le i\le P} \sum_{j\ne i} (p_i p_j)^{-\alpha}.
\]
\end{lemma}

\begin{proof}
Hermiticity is immediate from $D^*=D$ and $R^*=R$.
For the norm bound, we use $\|A_{\mathbb{P}}\| \le \|D\| + \|R\|$ and estimate each term separately.

Since $D$ is diagonal with real entries,
\[
\|D\| = \max_{1\le i\le P} |D_{ii}| = \max_{1\le i\le P} \sqrt{p_i}.
\]
To bound $\|R\|$, we use the inequality $\|R\| \le \max_i \sum_j |R_{ij}|$ (the $\ell^1$--row bound for the operator norm on $\ell^2$):
\[
\|R\|
  \le \max_{1\le i\le P} \sum_{j=1}^P |R_{ij}|
  =  \max_{1\le i\le P} \sum_{j\ne i} |\varepsilon_{ij}| (p_i p_j)^{-\alpha}
  \le \varepsilon_{\max} \max_{1\le i\le P} \sum_{j\ne i} (p_i p_j)^{-\alpha}.
\]
Combining these gives the desired bound.
\end{proof}

We lift $A_{\mathbb{P}}$ to $H_{P,N_t,d}$ by
\[
\mathbf{A} := A_{\mathbb{P}} \,\otimes\, I_{N_t} \,\otimes\, I_d.
\]

\begin{corollary}[Norm of $\mathbf{A}$]
$\mathbf{A}$ is Hermitian and
\[
\|\mathbf{A}\| = \|A_{\mathbb{P}}\|.
\]
\end{corollary}

\begin{proof}
Hermiticity follows from the tensor product of Hermitian operators.
For the norm, since $I_{N_t}$ and $I_d$ both have operator norm $1$, we have
\[
\|\mathbf{A}\| = \|A_{\mathbb{P}}\|\cdot \|I_{N_t}\|\cdot\|I_d\|
               = \|A_{\mathbb{P}}\|.
\]
\end{proof}

\subsection{Time block \texorpdfstring{$\mathbf{B}$}{B}: Fourier multiplier and bounds}

Let $\mathcal{T}=\{t_1,\dots,t_{N_t}\}$ be a uniform grid, and let $\{\omega_\ell\}_{\ell=0}^{N_t-1}$ denote the discrete Fourier frequencies.
Let $\mathsf{F}:\mathbb{C}^{N_t}\to\mathbb{C}^{N_t}$ be the discrete Fourier transform, normalized to be unitary.

Fix $\beta>0$ and define a real symbol
\[
m(\omega_\ell) := e^{-\beta \omega_\ell^2},\qquad \ell=0,\dots,N_t-1.
\]
Define $M$ on $\mathbb{C}^{N_t}$ by
\[
(M \hat f)_\ell = m(\omega_\ell)\hat f_\ell.
\]
Then $M$ is diagonal in the Fourier basis, with real diagonal entries $m(\omega_\ell)\in (0,1]$.
Define
\[
B_t := \mathsf{F}^{-1} M \,\mathsf{F}.
\]

\begin{lemma}[Boundedness and norm of $B_t$]
$B_t$ is a self-adjoint contraction on $\mathbb{C}^{N_t}$, with
\[
\|B_t\| = \max_{\ell} |m(\omega_\ell)| = 1.
\]
\end{lemma}

\begin{proof}
Since $\mathsf{F}$ is unitary and $M$ is a real diagonal matrix,
$B_t$ is unitarily equivalent to $M$, hence self-adjoint.
The operator norm of a normal operator equals the maximum modulus of its eigenvalues, so
\[
\|B_t\| = \|M\| = \max_\ell |m(\omega_\ell)| = 1,
\]
because $m(\omega_\ell)\in (0,1]$ and achieves $1$ at $\omega_\ell=0$.
\end{proof}

We extend $B_t$ to $H_{P,N_t,d}$ by
\[
\mathbf{B} := I_P \otimes B_t \otimes I_d.
\]

\begin{corollary}[Norm of $\mathbf{B}$]
$\mathbf{B}$ is Hermitian and
\[
\|\mathbf{B}\| = \|B_t\| = 1.
\]
\end{corollary}

\subsection{Internal block \texorpdfstring{$\mathbf{E}$}{E}}

Let $E\in M_d(\mathbb{C})$ be Hermitian, with eigenvalues $\lambda^{(E)}_1,\dots,\lambda^{(E)}_d$.
Define
\[
\mathbf{E} := I_P \otimes I_{N_t} \otimes E
\]
on $H_{P,N_t,d}$.

\begin{lemma}[Norm of $\mathbf{E}$]
$\mathbf{E}$ is Hermitian and
\[
\|\mathbf{E}\| = \|E\| = \max_{1\le k\le d} |\lambda^{(E)}_k|.
\]
\end{lemma}

\begin{proof}
Hermiticity is immediate.
The norm follows from the fact that tensoring with identities does not change the spectral radius, and for Hermitian operators the operator norm equals the spectral radius.
\end{proof}

\subsection{Universal operator \texorpdfstring{$\mathcal{U}$}{U}: basic bound}

We define the truncated universal operator by
\[
\mathcal{U} := \mathbf{A} + \mathbf{B} + \mathbf{E}.
\]

\begin{proposition}[Boundedness and norm estimate for $\mathcal{U}$]
$\mathcal{U}$ is Hermitian on $H_{P,N_t,d}$ and
\[
\|\mathcal{U}\|
  \le \|A_{\mathbb{P}}\| + \|B_t\| + \|E\|
  = \|A_{\mathbb{P}}\| + 1 + \|E\|.
\]
\end{proposition}

\begin{proof}
Sum of Hermitian operators is Hermitian.
The triangle inequality gives
\[
\|\mathcal{U}\| \le \|\mathbf{A}\| + \|\mathbf{B}\| + \|\mathbf{E}\|
                 = \|A_{\mathbb{P}}\| + 1 + \|E\|.
\]
\end{proof}

\section{K--theoretic multiplicity functors for $S^1$ and $T^2$}
\label{app:k-functors}

\subsection{$S^1$: rank functor}

Let $\mathcal{C}_{S^1}$ be the category of complex vector bundles over $S^1$, and recall that $K^0(S^1)\cong \mathbb{Z}$ via the rank.
Fix a finite set of \emph{base primes} $\mathcal{P}_{\mathrm{rank}} := \{p_{r,1},\dots,p_{r,q}\}\subset\mathcal{P}$.

For $E\in\mathcal{C}_{S^1}$ of rank $r\in\mathbb{Z}_{\ge 0}$, define its multiplicity vector
\[
v_E := \sum_{i=1}^q r\,\lvert p_{r,i}\rangle \in \ell^2(\mathcal{P}),
\]
with all other prime components zero.
Lift $v_E$ to a state in $H_{P,N_t,d}$ by tensoring with a fixed time state and internal state, e.g.
\[
\Psi_E^{(0)} := v_E \otimes \lvert t_1\rangle \otimes \lvert \alpha_0\rangle.
\]

\begin{proposition}[Functoriality for $S^1$]
The assignment
\[
\mathcal{M}_K^{(S^1)} : K^0(S^1)\to \ell^2(\mathcal{P}),\qquad r\mapsto v_E
\]
is a group homomorphism.
Moreover, for bundle maps $f:E\to F$ that preserve rank, the induced map on multiplicity vectors is the identity on $v_E$.
\end{proposition}

\begin{proof}
For $E,F$ of ranks $r,s$ respectively, the direct sum $E\oplus F$ has rank $r+s$, thus
\[
\mathcal{M}_K^{(S^1)}([E\oplus F])
  = \sum_{i=1}^q (r+s)\, \lvert p_{r,i}\rangle
  = \sum_{i=1}^q r\,\lvert p_{r,i}\rangle + \sum_{i=1}^q s\,\lvert p_{r,i}\rangle
  = \mathcal{M}_K^{(S^1)}([E]) + \mathcal{M}_K^{(S^1)}([F]),
\]
so the map is additive.
If $f:E\to F$ is a bundle map and $\operatorname{rank}(E)=\operatorname{rank}(F)$, then $[E]=[F]$ in $K^0(S^1)$, hence $\mathcal{M}_K^{(S^1)}([E])=\mathcal{M}_K^{(S^1)}([F])$.
\end{proof}

\subsection{$T^2$: rank and first Chern class}

For the toy model considered in the paper, we take $K^0(T^2)\cong \mathbb{Z}\oplus\mathbb{Z}$, parameterized by rank $r$ and an integer $c$ representing the first Chern class (in a chosen basis).\footnote{This is an effective parametrization in the torsion–free setting of interest.}
Fix two disjoint finite sets of primes
\[
\mathcal{P}_{\mathrm{rank}} := \{p_{r,1},\dots,p_{r,q}\},\qquad
\mathcal{P}_{\mathrm{chern}} := \{p_{c,1},\dots,p_{c,q'}\}\subset\mathcal{P},
\]
with $\mathcal{P}_{\mathrm{rank}}\cap\mathcal{P}_{\mathrm{chern}}=\varnothing$.

For $E_{r,c}$ with parameters $(r,c)\in\mathbb{Z}\oplus \mathbb{Z}$, define
\[
v_{r,c}
  := r \sum_{i=1}^q \lvert p_{r,i}\rangle
   + c \sum_{j=1}^{q'} \lvert p_{c,j}\rangle \in \ell^2(\mathcal{P}).
\]

\begin{proposition}[Additivity for $T^2$]
The map
\[
\mathcal{M}_K^{(T^2)}: K^0(T^2)\to \ell^2(\mathcal{P}),\qquad
(r,c)\mapsto v_{r,c}
\]
is a group homomorphism when $K^0(T^2)$ is regarded as $\mathbb{Z}\oplus\mathbb{Z}$ with the componentwise addition.
\end{proposition}

\begin{proof}
For $(r_1,c_1)$ and $(r_2,c_2)$,
\[
v_{r_1,c_1} + v_{r_2,c_2}
 = (r_1+r_2) \sum_i \lvert p_{r,i}\rangle
 + (c_1+c_2) \sum_j \lvert p_{c,j}\rangle
 = v_{r_1+r_2,\,c_1+c_2} = \mathcal{M}_K^{(T^2)}((r_1,c_1)+(r_2,c_2)).
\]
\end{proof}

\begin{remark}[Torsion–free regime]
For the $S^1$ and $T^2$ models used here, the relevant $K$--groups are torsion–free (isomorphic to $\mathbb{Z}$ or $\mathbb{Z}\oplus\mathbb{Z}$).
The multiplicity encodings above are therefore linear injections of these free abelian groups.
Extensions to spaces with torsion components will require a separate encoding scheme for finite cyclic factors and are left for future work.
\end{remark}

\section{Bridge operators and basic norm bounds}
\label{app:bridge}

\subsection{Finite zeta–bridge operator}

Let $\{\rho_k\}_{k=1}^M$ be a finite set of nontrivial zeros of the Riemann zeta function, $\rho_k = \tfrac12 + i\gamma_k$ with real $\gamma_k$.
For $p_i\in\mathcal{P}$, define phase factors
\[
\phi_{i,k} := \gamma_k \log p_i.
\]
Let $w_{i,k}\in\mathbb{C}$ be fixed coupling weights and $\{g_k(t_j)\}$ be scalar envelope factors (e.g.\ discrete Gaussians in $t_j$).
We define the truncated bridge operator
\[
\mathcal{B}_{P,M} : H_{P,N_t,d}\to H_{P,N_t,d}
\]
by its action on basis vectors
\[
\mathcal{B}_{P,M}\,\lvert i,j,\alpha\rangle
  := \sum_{k=1}^M w_{i,k}\, e^{\,i\phi_{i,k}}\,
     g_k(t_j) \,
     \lvert i,j,\alpha\rangle,
\]
so that $\mathcal{B}_{P,M}$ is diagonal in the $(i,j,\alpha)$ basis.

\begin{lemma}[Boundedness of the truncated bridge]
Assume $\sup_{i,k} |w_{i,k}| \le W_{\max}$ and $\sup_{j,k} |g_k(t_j)| \le G_{\max}$.
Then $\mathcal{B}_{P,M}$ is bounded with
\[
\|\mathcal{B}_{P,M}\|
  \le M\, W_{\max} G_{\max}.
\]
\end{lemma}

\begin{proof}
For each basis vector $\lvert i,j,\alpha\rangle$ we have
\[
\bigl\|\mathcal{B}_{P,M}\,\lvert i,j,\alpha\rangle\bigr\|
  = \biggl|\sum_{k=1}^M w_{i,k} e^{i\phi_{i,k}} g_k(t_j)\biggr|
  \le \sum_{k=1}^M |w_{i,k}|\, |g_k(t_j)|
  \le M\,W_{\max} G_{\max}.
\]
Since $\mathcal{B}_{P,M}$ is diagonal in the orthonormal basis $\{\lvert i,j,\alpha\rangle\}$, its operator norm is the supremum of the absolute diagonal entries, hence bounded by $M W_{\max} G_{\max}$.
\end{proof}

\subsection{Hermitian vs.\ unitary variants}

The above $\mathcal{B}_{P,M}$ is diagonal but not necessarily Hermitian.
For a Hermitian bridge one can symmetrize by choosing real $w_{i,k}$ and $g_k(t_j)$ and replacing $e^{i\phi_{i,k}}$ by $\cos(\phi_{i,k})$, so that $\mathcal{B}_{P,M}$ is real symmetric.

For a unitary bridge, one instead defines a diagonal unitary
\[
U_{\mathcal{B}}\,\lvert i,j,\alpha\rangle
  := \exp\!\Bigl(i \sum_{k=1}^M \theta_{i,j,k}\Bigr)\lvert i,j,\alpha\rangle,
\]
with real phases $\theta_{i,j,k}$ chosen so that the combined phase
$\sum_k \theta_{i,j,k}$ approximates $\sum_k w_{i,k} g_k(t_j)\phi_{i,k}$ in a suitable sense.
Then $\|U_{\mathcal{B}}\|=1$ exactly.

\section{Norm bound for the coupled operator}
\label{app:coupled}

Consider the coupled operator
\[
\mathcal{U}_{\mathrm{cpl}} := \mathcal{U} + \lambda\,\mathcal{B}_{P,M},
\]
with $\lambda\in\mathbb{R}$.
Under the boundedness assumptions above we obtain:

\begin{proposition}[Norm estimate for $\mathcal{U}_{\mathrm{cpl}}$]
Assume $\mathcal{B}_{P,M}$ satisfies $\|\mathcal{B}_{P,M}\|\le C_{P,M}$.
Then
\[
\|\mathcal{U}_{\mathrm{cpl}}\|
  \le \|A_{\mathbb{P}}\| + 1 + \|E\| + |\lambda| C_{P,M}.
\]
\end{proposition}

\begin{proof}
By the triangle inequality,
\[
\|\mathcal{U}_{\mathrm{cpl}}\|
  \le \|\mathcal{U}\| + |\lambda|\|\mathcal{B}_{P,M}\|
  \le (\|A_{\mathbb{P}}\| + 1 + \|E\|) + |\lambda| C_{P,M}.
\]
\end{proof}

This bound shows that, for fixed $P,M,N_t$ and bounded weights $(w_{i,k},g_k)$, the truncated ZMOS operators are uniformly bounded in operator norm.

\section{Finite--truncation stability considerations}
\label{app:stability}

Let $(P,M,N_t)$ denote truncation parameters and write $\mathcal{U}^{(P,M,N_t)}_{\mathrm{cpl}}$ for the corresponding coupled operator.
Let $\{\lambda_i^{(P,M,N_t)}\}$ denote its eigenvalues (counted with multiplicity).

\begin{definition}[Finite--truncation stability]
A spectral statistic $S$ (e.g.\ a normalized gap distribution, an averaged observable, or the dynamical entropy proxy $\Delta(t)$) is said to be \emph{stable under truncation} if for every $\varepsilon>0$ there exist $P_0,M_0,N_{t,0}$ such that for all $P\ge P_0$, $M\ge M_0$, $N_t\ge N_{t,0}$,
\[
\bigl| S(\mathcal{U}^{(P,M,N_t)}_{\mathrm{cpl}}) - S_\infty \bigr| < \varepsilon,
\]
for some limiting value $S_\infty$ (which may itself depend on other model parameters).
\end{definition}

In the present work we do not prove such stability analytically; instead, we adopt a \emph{finite--truncation--first} philosophy:
\begin{itemize}
  \item All spectral and dynamical claims are first formulated for fixed $(P,M,N_t)$.
  \item Numerical experiments are used to check that the corresponding statistics exhibit convergence as $(P,M,N_t)$ increase within computationally accessible ranges.
\end{itemize}
A rigorous stability analysis for specific operator families is left for future work.

\bigskip

The results in this appendix give explicit constructions and norm bounds for the building blocks of ZMOS in finite truncation.
They guarantee that the operators considered in the main text are bounded (and Hermitian or unitary as required), and they formalize the $K$--theoretic multiplicity encodings used for $S^1$ and $T^2$ in the torsion--free regime.


\nocite{*}
\bibliographystyle{plainnat}
\bibliography{references}
\end{document}